% === Begin Label Data ===
% ID: 879
% 难度星级: 
% 标签: 
% 备注: 新定义
% 组卷引用次数: 0
% === End  Label Data ===

\begin{problem}{2021}{G}{北京卷}{21}{数列}
设 $p$ 为实数. 若无穷数列 $\{a_n\}$ 满足如下三个性质,则称 $\{a_n\}$ 为 $\mathfrak{R}_p$ 数列:
\circled{1}$ a_1 + p \geqslant 0$,且 $a_2 + p = 0$;
\circled{2}$ a_{4n-1} < a_{4n} (n=1,2,\dots)$;
\circled{3}$ a_{m+n} \in \{a_m + a_n + p, a_m + a_n + p + 1\} (m=1,2,\dots; n=1,2,\dots)$.

( I )如果数列 $\{a_n\}$ 的前四项为 $2,-2,-2,-1$,那么 $\{a_n\}$ 是否可能为 $\mathfrak{R}_2$ 数列?说明理由;
( II )若数列 $\{a_n\}$ 是 $\mathfrak{R}_0$ 数列,求 $a_5$;
( III)设数列 $\{a_n\}$ 的前 $n$ 项和为 $S_n$,是否存在 $\mathfrak{R}_p$ 数列 $\{a_n\}$,使得 $S_n \geqslant S_{10}$ 恒成立?如果存在,求出所有的 $p$;如果不存在,说明理由.
\end{problem}

\begin{answer}
( I )不可能; ( II ) $a_5=1$; ( III)存在, $p=2$.
\end{answer}

\begin{solutions}
解:( I )数列 $\{a_n\}$ 不可能为 $\mathfrak{R}_2$ 数列,理由如下,

因为 $p=2, a_1=2, a_2=-2$,所以 $a_1+a_2+p=2, a_1+a_2+p+1=3$,
因为 $a_3=-2$,所以 $a_3 \notin \{a_1+a_2+p, a_1+a_2+p+1\}$,
所以数列 $\{a_n\}$ 不满足性质 \circled{3}.

( II )性质 \circled{1}$ a_1 \geqslant 0, a_2=0$;
由性质\circled{3} $ a_2 \in \{a_1+a_1+0, a_1+a_1+0+1\}$,
则 $2a_1=0$ 或 $2a_1+1=0$,后者与 $a_1 \geqslant 0$ 矛盾,故 $a_1=0$;
由性质 \circled{3}$ a_{m+2} \in \{a_m, a_m+1\}$,因此 $a_3=0$ 或 $a_3=1, a_4=0$ 或 $a_4=1$,
通过性质 \circled{2} 只能是 $a_4=1, a_3=0$,
由性质 \circled{3}$ a_5=a_{2+3} \in \{0,1\}$ 或 $a_5=a_{1+4} \in \{1,2\}$
综上 $a_5=1$.

(III)先探求必要条件. 满足题设要求的 $R_p$ 数列 $\{a_n\}$,其前 $n$ 项和 $S_n$ 的最小值是 $S_{10}$,
所以有必要条件 $a_{10} \leqslant 0 \leqslant a_{11}$.
$-p=a_2 \in \{2a_1+p, 2a_1+p+1\} \Leftrightarrow a_1=-p$.
因为 $a_3 \in \{-p, -p+1\}, a_4 \in \{-p, -p+1\}$,
结合 $a_3 < a_4$,所以 $a_3=-p, a_4=-p+1$.
因为 $a_5 \in \{-p+1, -p+2\} \cap \{-p, -p+1\} = \{-p+1\}$,所以 $a_5=-p+1$.
因为 $a_6 \in \{-p+1, -p+2\} \cap \{-p, -p+1\} = \{-p+1\}$,所以 $a_6=-p+1$.
因为 $a_7 \in \{-p+1, -p+2\}, a_8 \in \{-p+2, -p+3\}$,
结合 $a_7 < a_8$,所以 $a_7=-p+1, a_8=-p+2$.
因为 $a_9 \in \{-p+2, -p+3\} \cap \{-p+1, -p+2\} = \{-p+2\}$,所以 $a_9=-p+2$.
因为 $a_{10} \in \{-p+2, -p+3\} \cap \{-p+1, -p+2\} = \{-p+2\}$,所以 $a_{10}=-p+2$.
因为 $a_{11} \in \{-p+2, -p+3\}, a_{12} \in \{-p+3, -p+4\}$
结合 $a_{11} < a_{12}$,所以 $a_{11}=-p+2, a_{12}=-p+3$.
故由必要条件得 $-p+2 \leqslant 0 \leqslant -p+2$,故 $p=2$.
再证充分性:即证 $R_2$ 数列的前 $n$ 项和最小值是 $S_{10}$.

方法一——猜想验证成立:
$a_n=\begin{cases} -2+k, n \in \{4k+1, 4k+2, 4k+3\} \\ -1+k, n=4k+4 \end{cases} (k \in \mathbb{N})$.
下面验证数列 $\{a_n\}$ 满足性质(1)(2)(3).
由 $a_1=-2, a_2=-2$ 可知 $a_1+p \geqslant 0, a_2+p=0$.
因为 $a_{4n-1}=n-3, a_{4n}=n-2$,所以 $a_{4n-1} < a_{4n}$.
对于任意正整数 $m,n$,存在 $k_1, k_2 \in \mathbb{N}, r_1, r_2 \in \{0,1,2,3\}$,使得 $m=4k_1+r_1, n=4k_2+r_2$.
所以 $a_m=-2+k_1, a_n=-2+k_2$.
所以 $a_m+a_n+p=-2+k_1+k_2$,
$a_m+a_n+p+1=-1+k_1+k_2$.
又 $m+n=4(k_1+k_2)+r_1+r_2$,所以
当 $0 \leqslant r_1+r_2 < 4$ 时, $a_{m+n}=-2+k_1+k_2$;
当 $4 \leqslant r_1+r_2 < 6$ 时, $a_{m+n}=-1+k_1+k_2$.
所以 $a_{m+n} \in \{a_m+a_n+p, a_m+a_n+p+1\}$.
由通项公式可知,当 $n < 9$ 时, $a_n \leqslant a_{10}=0$;
当 $n > 11$ 时, $a_n \geqslant a_{10}=0$.
所以 $S_n \geqslant S_{10}$ 恒成立.
综上,存在 $\mathfrak{R}_p$ 数列 $\{a_n\}$ 使得 $S_n \geqslant S_{10}$ 恒成立,这时 $p=2$.

方法二——数学归纳法证明:
下面用数学归纳法证明 $a_{4n+i}=-p+n (i=1, 2, 3), a_{4n+4}=-p+n+1 (n \in \mathbb{N})$,

当 $n=0$ 时,经检验命题成立;

假设 $n=k (k \geqslant 0)$ 时命题成立.

当 $n=k+1$ 时,

若 $i=1$,则 $a_{4(k+1)+1}=a_{4k+5}=a_{1+(4k+4)} \in \{-p+k+1, -p+k+2\}=a_{4k+5}=a_{2+(4k+3)} \in \{-p+k, -p+k+1\}$

综上 $a_{4k+5}=-p+k+1$;

若 $i=2$, $a_{4(k+1)+2}=a_{4k+6}=a_{2+(4k+4)} \in \{-p+k+1, -p+k+2\}$

$=a_{4k+6}=a_{3+(4k+3)} \in \{-p+k, -p+k+1\}$

综上 $a_{4k+6}=-p+k+1$,

若 $i=3$, $a_{4(k+1)+3}=a_{4k+7}=a_{3+(4k+4)} \in \{-p+k+1, -p+k+2\}$

结合 $a_{4(k+1)+4}=a_{4k+8}=a_{3+(4k+5)} \in \{-p+k+1, -p+k+2\}$

通过性质 \circled{3} 得 $a_{4k+7}=-p+k+1, a_{4k+8}=-p+k+2$,

即当 $n=k+1$ 时,命题成立.

数列单调,前 $n$ 项和最小值是 $S_{10}$,证毕.

方法三——研究单调性:
由条件 \circled{3} 得 $a_{n+1}=a_n+a_1+p$ 或 $a_{n+1}=a_n+a_1+p+1$.

又由条件 \circled{1} 得 $a_{n+1} \geqslant a_n$,则数列为单调递增数列,

$R_2$ 数列的前 $n$ 项和最小值是 $S_{10}$,证毕.
\end{solutions}